\section{Diffusion Models}
Another class of explicit density estimation generative models is \textbf{diffusion models}. 
Although initially proposed in 2015, they have only recently surged in popularity due to their ability to generate exceptionally high-quality samples. 
Around 2015, Generative Adversarial Networks (GANs) were the state-of-the-art in generative modeling, despite having some serious limitations (as discussed in the previous chapter). 
It was not until 2020 that breakthroughs in architectures and training techniques allowed diffusion models to surpass GANs in sample quality, establishing them as the current state-of-the-art.

Similar to VAEs, diffusion models use neural networks to transform a simple prior distribution (e.g., a standard Gaussian) into a complex data distribution that ideally matches the training data.

At a high level, the fundamental goal of this generative framework involves two simple steps:
\begin{enumerate}
    \item Sample a latent variable $z \sim \mathcal{N}(0, I)$ from a simple distribution;
    \item Pass $z$ through a neural network to generate a sample $\hat{x}$.
\end{enumerate}

While traditional models like GANs might do this in a single pass, diffusion models achieve this transformation iteratively, gradually denoising the latent sample step-by-step. Ultimately, the network is trained to minimize the discrepancy between the generated data and the real data distribution.

Two key questions arise from this overarching framework:
\begin{itemize}
    \item How do we define and structure the latent space distribution?
    \item How do we measure distribution similarity (i.e., what objective function do we use to train the neural network)?
\end{itemize}

\subsection{Idea}
The core idea behind diffusion models is actually quite simple: 
we start with a sample from the data distribution and gradually add noise to it over a series of steps, until it becomes pure noise. 
Then, we train a neural network to reverse this process, learning to denoise the noisy samples step-by-step until we recover the original data sample.

More formally, we define:
\begin{itemize}
    \item \textbf{Forward diffusion process:} a Markov chain that gradually adds noise to the data sample $x_0$ over $T$ steps, 
    resulting in a sequence of noisy samples $x_1, x_2, \ldots, x_T$, where $x_T$ is nearly pure noise.
    \item \textbf{Reverse denoising process:} a neural network that learns to reverse the forward diffusion process,
    starting from the noisy sample $x_T$ and iteratively denoising it to recover the original data sample $x_0$.
\end{itemize}

Let's now dive into the details of the forward diffusion process and the reverse denoising process, and how they are implemented in practice.

\subsection{Forward Process}
The forward diffusion process can be defined as a fixed Markov chain that gradually adds Gaussian noise to the data sample $x_0$ over $T$ steps.

\termdef{Markov chain}
{
    A Markov chain is a stochastic process that satisfies the \textbf{Markov property}: 
    the probability distribution of the next state depends only on the current state and not on the sequence of previous states. 
    This property is also known as \textbf{memorylessness}.
}

This implies that the forward process requires no learning, since it is completely specified by a predefined noise schedule.

We can formalize the forward diffusion process as follows:
\begin{equation}
    \label{eq:forward-diffusion}
    q(x_t | x_{t-1}) = \mathcal{N}(\sqrt{{1 - \beta_t}} x_{t-1}, \beta_t I)
\end{equation}
where $\beta_t$ is a variance schedule that controls the amount of noise added at each step. 

Basically, at each step $t$, we take the previous sample $x_{t-1}$, reduce its magnitude by a factor of $\sqrt{1 - \beta_t}$, subsituting it with a random Gaussian noise with variance $\beta_t$.

\subsubsection{A problem}
Suppose we have a data sample $x_0$ and we want to generate its noisy version $x_T$ at step $T$.
We can do this by iteratively applying the forward diffusion (Equation \ref{eq:forward-diffusion}) process $T$ times, starting from $x_0$ and generating $x_1, x_2, \ldots, x_T$.
This would result in:
\[
    q(x_T | x_0) = \prod_{i=1}^{T} q(x_i | x_{i-1})
\]
However, this approach is computationally expensive, especially for large $T$, since we would need to sample from the Gaussian distribution $T$ times, which can be slow and inefficient.

\subsubsection{Reparametrization trick}
Instead, we can use a reparametrization trick to directly sample $x_T$ from $x_0$ in a single step, without having to go through all the intermediate steps.

First, we can rewrite the forward diffusion process in a more convenient form:
\[
    q(x_t | x_{t-1}) = \mathcal{N}(\sqrt{\alpha_t} x_0, (1 - \alpha_t) I)
\]
where $\alpha_t = (1 - \beta_t)$.

By recursively applying this equation, we can derive a closed-form expression for $q(x_T | x_0)$:
\[
    q(x_T | x_0) = \mathcal{N}(\sqrt{\bar{\alpha}_T} x_0, (1 - \bar{\alpha}_T) I)
\]
where $\bar{\alpha}_T = \prod_{t=1}^{T} \alpha_t$.

This makes training and sampling actually feasible, since we can now sample $x_T$ directly from $x_0$ in a single step, without having to go through all the intermediate steps.

\subsection{Reverse Process}
The reverse denoising process is the core of the diffusion model, where we train a neural network to learn the reverse of the forward diffusion process. 
The goal is to start from a noisy sample $x_T$ and iteratively denoise it to recover the original data sample $x_0$.

We could define the reverse process as a Markov chain that learns to reverse the forward diffusion process $p(x_{t-1} | x_t)$.
Applying Bayes' theorem, we can express the reverse process as:
\[
    p(x_{t-1} | x_t) = \frac{p(x_t | x_{t-1}) p(x_{t-1})}{p(x_t)}
\]

But this is intractable in practice, since the marginal distribution $p(x_t)$ is unknown and difficult to compute.

Therefore, as previously seen in the VAE chapter, we can train a neural network to approximate the reverse process $p_\theta(x_{t-1} | x_t)$, where $\theta$ are the parameters of the neural network.
By assuming $T \rightarrow \infty$ and $\beta_t \rightarrow 0$, we can show that the reverse process can be approximated as a Gaussian distribution:
\[
    p_\theta(x_{t-1} | x_t) = \mathcal{N}(\mu_\theta(x_t, t), \Sigma_t)
\]
where $\mu_\theta(x_t, t)$ is the mean predicted by the neural network and $\Sigma_t$ is a fixed variance schedule.

Usually, only the mean $\mu_\theta(x_t, t)$ is predicted by the neural network, while the variance $\Sigma_t$ is fixed and known. 
It was later shown that learning the variance $\Sigma_t$ can improve sample quality, but it is not strictly necessary.

\subsubsection{Hierchical VAEs}
We can also view diffusion models as a special case of hierarchical VAEs, with some important properties:
\begin{itemize}
    \item \textbf{Fixed Encoder:} the encoder is fixed and defined by the forward diffusion process;
    \item \textbf{Latent Variable:} all the latent variables $x_1, x_2, \ldots, x_T$ have the same dimensionality as the data sample $x_0$;
    \item \textbf{Reusable Decoder:} the neural network decoder is reused at each step of the reverse process, conditioned on the current step $t$;
\end{itemize}

